\ssr{ВВЕДЕНИЕ}

Цель работы -- сравнить последовательную и конвейерную обработку.

Задачи, необходимые для достижения поставленной цели:
\begin{enumerate}
	\item описать и реализовать задание по варианту;
	\item сравнить реальное время, затрачиваемое на обработку N заявок последовательно и конвейером.
\end{enumerate}

% Конвейер, согласно общему заданию: k потоков, включая 3 отдельных на каждое ОУ и (k-3) дополнительно для ОУ2. k взять как рекомендованное значение для выбранной АЭВМ из ЛР4.

% Обслуживающее устройство 1 выполняет считывание графа в формате матрицы смежности из текстового или бинарного файла (имя файла дано в заявке). Обслуживающее устройство 2 выполняет обработку графа, согласно заданию на ЛР4, задействовав k-3 вспомогательных потоков. Обслуживающее устройство 3 выполняет дамп ответа в файл, имя которого получено добавлением префикса к имени входного файла из заявки (файлы с ответами соотносятся с заявками 1:1).

% 3 нативных потока, которые передают друг другу информацию через очереди. Данные в 1-ю очередь (входную для 1-го обслуживающего устройства) в количестве N заявок заполнить перед пуском конвейера, считав их из файла (список имен файлов).

% Сравнить реальное время, затрачиваемое на обработку N \in {25, 50, 75, 100, 125} заявок последовательно и конвейером, основанным на k нативных потоках, не считая диспетчера.

% Пример сформированного программой лога (см. в почте фотокопии материалов лекции) поместить после тестирования, чтобы пронаблюдать, что некоторые события происходят (почти) одновременно, и, следовательно, заявки действительно обслуживаются параллельно.

% Ограничения/допущения --- см. в почте фотокопии материалов лекции.

\clearpage
